\section{Conclusions}
\label{sec:conclusions}

This work provides a reproducible method for fitting power diagrams to selected numerical separator observations at fixed sites. The affine relation between separator positions and weight differences reduces exact compatibility and identifiability to graph conditions, and unconstrained quadratic fitting to a weighted graph-Laplacian problem. Compatible observations determine relative weights within each connected component; a connected observation graph determines all weights up to a common shift that leaves the diagram unchanged. Synthetic tests with exact compatible data recover these observable differences to numerical precision and show how missing observations limit reconstruction.

Geometric assessment asks whether the cells share a face and whether it contains the fitted plane's intersection with the line through the sites. When this point lies inside the face, centered clearance measures its in-plane distance to the boundary. For finite nonperiodic configurations, a strict power-distance advantage of the pair over all competing sites at this point certifies a face nearby, and the perturbation bounds give sufficient conditions for retaining that local advantage. Clipping and periodicity require the corresponding domain and image checks.

The experiments show that agreement with supplied separators and with an independent reference can diverge. In synthetic tests, better target fitting can accompany poorer recovery of reference adjacencies. In the 100-molecule application, fits using the known reference pairs realize all their requested faces, although reconciling local predictions from atomic density profiles does not necessarily improve separator accuracy against electron-density landmarks. With distance-selected candidates, connector-visibility filtering retains 3,973 of the 3,975 reference pairs, together with additional nonbonded pairs. Centered clearance adds further information: among visible nonbonded faces, it ranks density-supported reference pairs more effectively than distance or face area within each model's evaluated population. These findings concern local separator positions and pair relations; complete electron-density basins have not been reconstructed or validated.

Optional face-driven refinement improves accuracy in some tests and degrades it in others, and its stopping diagnostics do not establish general convergence. The main outcome is a method for turning incomplete local separator information into a global power diagram whose identifiable weight information, realized faces and local geometric support can each be examined explicitly.

